% ===================================================================
%  نقشہ نمبر 11 — شہر تے اوہدیاں عمارتاں۔ — اگ۔
%  Traduction 2026 d'apres texto/io, controlee sur texto/fr, texto/en
%  et texto/hi. Consigne : voir texto/pnb/10-naqsha-01.tex.
% ===================================================================

%%K t11-apar-1 apar
\VUsaut{2.0mm}
\VUcentre{13.2pt}{90}{تیجا سلسلہ}
\VUsaut{3.0mm}
\VUfilet{16mm}
\VUsaut{5.6mm}
\VUtitre{15.0pt}{60}{نقشہ نمبر 11}
\VUsaut{1.4mm}
\VUpk{11.4pt}{40}{شہر تے اوہدیاں عمارتاں۔ --- اگ۔}
\VUsaut{2.2mm}

%%K t11-tit-1 sub
\VUpk{10.2pt}{40}{باہرلی بستی۔}
\VUsaut{3.0mm}

%%K t11-c1-01-1 p
1. --- میں محیط توں سو کلومیٹر دور اک وڈے شہر وچ رہنا واں، جیہڑا فیر وی پہلے
درجے دی \VUgras{بندرگاہ} \textsuperscript{(1)} اے۔ جیہڑا دریا اوہدے نال وگدا اے
اوہ چار سو میٹر چوڑا اے تے اک بہت سوہݨا لنگر گاہ بݨاندا اے۔

%%K t11-c1-02-1 p
2. --- \VUgras{گھاٹاں} \textsuperscript{(2)} نال دا سارا حصہ بالکل سمندری تے
صنعتی رُوپ رکھدا اے، تے اک ایہو جیہی \VUgras{باہرلی بستی} \textsuperscript{(3)}
بݨاندا اے جیدے وچ نِرے ملاح تے مزدور رہندے نیں۔ ایہہ \VUgras{کارخانیاں}
\textsuperscript{(4)} وچ تے \VUgras{گیس گھر} \textsuperscript{(5)} وچ کم کردے نیں،
جیدی اُچی \VUgras{چمنی} \textsuperscript{(6)} تے \VUgras{گیس ٹینکی} دور توں ای
وکھائی دیندی اے۔

%%K t11-c1-03-1 p
3. --- وچکارلے جُگ وچ شہر قلعہ بند سی، تے اوہدیاں فصیلاں دے کجھ نشان ہُݨ وی لبھدے
نیں، خاص کر کے اوہ \VUgras{بُوہا} \textsuperscript{(8)} جیہڑا پرانے
\VUgras{گڑھ قلعے} \textsuperscript{(7)} دا اے، تے جیدے دوہاں پاسیں اوہدے دو
\VUgras{منارے} \textsuperscript{(9)} نیں، جیہڑے ہُݨ \VUgras{قید خانے} دے کم آندے نیں۔

%%K t11-c1-04-1 p
4. --- ایس سارے \VUgras{محلے} وچ، جیہڑا بہت پرانا لگدا اے، نِری اک عمارت رتی نویں
اے: \VUgras{چنگی گھر} \textsuperscript{(10)}۔ اوہ گھاٹ تے اوس نِکی \VUgras{گلی}
\textsuperscript{(11)} وچکار کھڑا اے جیہڑی پرانے \VUgras{شہری گھر}
\textsuperscript{(12)} تیکر جاندی اے۔ اوس چھیکرلی عمارت نوں اوس وڈے
\VUgras{گھنٹہ منارے} \textsuperscript{(13)} توں سوکھا ای پچھاݨیا جا سکدا اے جیہڑا
پرانے شہر اُتے چھایا رہندا اے۔

%%K t11-c1-05-1 p
5. --- گلی دے دوجے پاسے چنگی گھر دی ای وضع دی اک عمارت \VUgras{کھڑی} اے: اوہ
\VUgras{وپار گھر} \textsuperscript{(14)} اے۔ اوہدا اک مونہہ اک نِکے چونک ول کھلدا اے
جِتھے \VUgras{انتظامی گھر} \textsuperscript{(15)} اے۔ ساڈا شہر، اصل وچ، راجدھانی نہ
ہوندے ہوئے وی اک اہم ضلعی صدر مقام اے۔

%%K t11-tit-2 sub
\VUsaut{5.0mm}
\VUpk{10.2pt}{40}{شہر دا اُتلا حصہ۔}
\VUsaut{3.4mm}

%%K t11-c1-06-1 p
6. --- چھاؤݨی، جیہڑی کافی وڈی اے، وڈیاں تے بہت صحت مند \VUgras{بیرکاں}
\textsuperscript{(16)} وچ رہندی اے؛ اوہ \VUgras{شہری باغ} \textsuperscript{(17)} دے
جنگلے تیکر پھیلیاں نیں۔ جیہڑی \VUgras{چوڑی سڑک} \textsuperscript{(19)} ایس باغ تیکر
جاندی اے، اوہ \VUgras{بچت بینک} \textsuperscript{(18)} دے پچھوں لنگھ کے
\VUgras{وڈے گرجے} \textsuperscript{(20)} دے چونک اُتے نکلدی اے۔

%%K t11-c1-07-1 p
7. --- ایہہ ساریاں عمارتاں اک نِکی ٹبی اُتے پوڑیاں وانگوں چڑھدیاں نیں، جِتھے ساڈا
وڈا \VUgras{ثانوی سکول} \textsuperscript{(21)} وی اے۔

%%K t11-c1-07-2 p
جے تُسیں اوس ہولی ڈھال والی سڑک توں اُترو جیہڑی وڈی سڑک نال ملدی اے، تے تُسیں
\VUgras{عدالت} \textsuperscript{(22)} تیکر اپڑدے او، جیدے سامݨے اک سوہݨا
\VUgras{عوامی باغ} \textsuperscript{(26)} اے۔ ایہہ \VUgras{چونک} بہت وڈا نئیں، پر شہر
دے وچکار اوہ ہریاول تے ٹھنڈ دا نخلستان بݨاندا اے۔ سو شہر والے اوہدے وچ خوب آندے
نیں، جیہناں نوں \VUgras{کہوہ خانے} \textsuperscript{(25)} دے چھجے دے تھلے بہہ کے اوس
فوجی باجے نوں سُݨنا چنگا لگدا اے جیہڑا ویروار تے اتوار نوں لان تے کیاریاں وچکار
بݨے \VUgras{باجے دے منچ} \textsuperscript{(27)} وچ وجدا اے۔

%%K t11-c1-08-1 p
8. --- اوہناں ای لاناں وچوں اک اُتے کانسی دی اک \VUgras{مورتی} \textsuperscript{(28)}
کھڑی اے، ساڈے اک شہر والے دی یاد وچ چُکی یادگار، جیہنے اپݨی جنم بھوں دی وڈی سیوا
کیتی۔

%%K t11-tit-3 sub
\VUsaut{4.0mm}
\VUpk{10.2pt}{40}{آوݨ جاوݨ۔}
\VUsaut{3.4mm}

%%K t11-c1-09-1 p
9. --- ساڈا شہر ہالے بجلی نال روشن نئیں؛ اوہدے وچ نِرے \VUgras{گیس دے لیمپ}
\textsuperscript{(29)} نیں۔ پر \VUgras{مقامی اخبار} ایس گل لئی مہم چلا رہے نیں پئی
روشنی دا اوہ بندوبست سُدھاریا جاوے، جِویں ٹراماں دا \VUgras{کھِچݨ دا بندوبست}
پہلاں ای سُدھاریا جا چُکیا اے۔ \VUgras{گھوڑیاں} نال \VUgras{کھِچیاں} جاوݨ والیاں
پرانیاں گڈیاں دی تھاں ہُݨ کم دیاں \VUgras{بجلی دیاں ٹراماں} \textsuperscript{(31)}
آ گئیاں نیں، جیہڑیاں بہت \VUgras{گھٹ} کرائے وچ مسافراں نوں \VUgras{شہر دے} اک سرے
توں \VUgras{دوجے تیکر} چھیتی اپڑا دیندیاں نیں۔

%%K t11-c1-10-1 p
10. --- عدالت کول \VUgras{بینک} \textsuperscript{(32)} اے، جِتھے وپاری ہنڈیاں بھنائیاں
جاندیاں نیں۔ \VUgras{بینک دے ہرکارے} \textsuperscript{(33)} لوکاں دے گھر گھر جا کے
بقایا وصول کردے نیں۔

%%K t11-tit-4 sub
\VUsaut{4.0mm}
\VUpk{10.2pt}{40}{کلا تے ودیا۔}
\VUsaut{3.4mm}

%%K t11-c1-11-1 p
11. --- کول دی اوہ سوہݨی عمارت \VUgras{عجائب گھر} اے۔ اوہدے وچ اکو ویلے شہر دا
\VUgras{کتب خانہ} \textsuperscript{(34)}، کجھ سوہݨے
\VUgras{قدرتی تاریخ دے سنگرہ} \textsuperscript{(35)} (اک قدرتی تاریخ دا عجائب گھر)
تے اک سوہݨا \VUgras{تصویر گھر} \textsuperscript{(36)} اے، جیہڑا اک شاندار پکی
\VUgras{نمائش} بݨاندا اے۔

%%K t11-c1-11-2 p
اوہ ہفتے وچ دو واری لوکاں لئی کھلدا اے، پر ویکھݨ والے کِسے وی دن اوہنوں ویکھ سکدے
نیں، \VUgras{رکھوالے} \textsuperscript{(37)} نوں تھوڑی بخشش دے کے۔ \VUgras{مصور}
\textsuperscript{(38)} اوتھے کم کرن آندے نیں۔ اوہناں نوں اپݨیاں تپائیاں دے
\VUgras{سامݨے} ویکھیا جا سکدا اے، ہتھ وچ \VUgras{رنگ پٹی} چُکی، مشہور تصویراں دی
نقل اُتاردے۔

%%K t11-c1-12-1 p
12. --- عجائب گھر دے پچھلیاں اُچائیاں اُتے اوہ \VUgras{گنبد} \textsuperscript{(40)}
جیہڑا \VUgras{ہسپتال} \textsuperscript{(39)} دا اے، تے اوس
\VUgras{ماسٹراں دے تربیتی سکول} \textsuperscript{(41)} دیاں عمارتاں مشکل نال وکھائی
دیندیاں نیں جِتھے ساڈے شہری سکولاں دے ماسٹر تربیت پائے سن۔ رنگ منچ دے چونک اُتے اک
\VUgras{ڈاک خانہ} \textsuperscript{(43)} اے، جیہڑا اک \VUgras{نجی مکان}
\textsuperscript{(42)} وچ بیٹھا اے۔

%%K t11-tit-5 sub
\VUsaut{5.0mm}
\VUpk{11.4pt}{40}{اگ۔}
\VUsaut{3.0mm}

%%K t11-tit-6 sub
\VUpk{10.2pt}{40}{اگ دی ہاک۔}
\VUsaut{3.4mm}

%%K t11-c2-01-1 p
1. --- شہر وچ ڈراؤݨی خبر پھیل رہی اے: رنگ منچ نوں ہُݨے ہُݨے اگ لگ گئی! جدوں تیکر
میں \VUgras{چونک} \textsuperscript{(96)} اپڑنا واں، بھیڑ \VUgras{پٹڑی}
\textsuperscript{(46)} تے \VUgras{سڑک} \textsuperscript{(47)} اُتے جُڑ چُکی اے۔
\VUgras{ڈاک دے بابو} \textsuperscript{(44)} تے \VUgras{ڈاکیے} \textsuperscript{(45)}
ڈاک خانے توں باہر آ رہے نیں۔ اک \VUgras{ٹرام چالک} \textsuperscript{(48)} نوں اپݨی
ٹرام روکݨی پئی اے تاں جو شہر دی اک \VUgras{مریض گڈی} \textsuperscript{(50)} نکل سکے،
جیہڑی اک \VUgras{جراح} \textsuperscript{(52)} تے اک \VUgras{نرس} \textsuperscript{(51)}
نوں لے کے آ رہی اے، اک \VUgras{زخمی} \textsuperscript{(53)} دی سنبھال لئی، جیہنوں
\VUgras{اسٹریچر} \textsuperscript{(54)} اُتے \VUgras{اخبار دے کھوکھے}
\textsuperscript{(55)} کول لیایا گیا اے۔

%%K t11-c2-02-1 p
2. --- مشق توں مُڑدی پیدل فوج دی اک ٹولی رُک گئی اے، تاں جو گھوڑیاں اُتے سوار
\VUgras{شہری محافظاں} \textsuperscript{(61)} نال بندوبست بݨائی رکھݨ وچ مدد کرے۔
\VUgras{گھوڑ سواراں} دے گھوڑے دو پیراں اُتے کھڑے ہو جاندے نیں؛ اوہ \VUgras{سپاہیاں}
\textsuperscript{(57)} یا \VUgras{ژندارماں} \textsuperscript{(63)} نالوں چنگی طرحاں
\VUgras{تماشبیناں} \textsuperscript{(56)} نوں پچھے دھکّ سکدے نیں۔ ژندارم ویسے وی بہت
رُجھے نیں۔ اوہناں وچوں اک \VUgras{پولیس دے سپاہیاں} \textsuperscript{(64)} نوں دس
رہیا اے پئی اک ہور \VUgras{زخمی} \textsuperscript{(65)} نوں کِتھے لے جاوݨا اے، اک
سورما اگ بجھاؤ والا جیہڑا پوڑی توں ڈِگ پیا اے۔

%%K t11-c2-03-1 p
3. --- \VUgras{افسر} \textsuperscript{(66)} ٹھیک ٹھیک دسدا اے پئی جیہڑا
\VUgras{بھاپ دا پمپ} \textsuperscript{(67)} ہُݨے آ رہیا اے اوہنوں کِتھے رکھݨا اے۔ اوس
تکڑے جنتر دی اُڈیک وچ اوہنے اک \VUgras{ہتھ دا پمپ} \textsuperscript{(68)} لگوا دِتا
اے، جیدی لمی \VUgras{نالی} \textsuperscript{(69)} اگ اُتے پاݨی دیاں دھاراں ڈوہل رہی
اے۔ چھے اگ بجھاؤ والے اوہنوں چلا رہے نیں، جدوں تیکر اوہناں دا ساتھی \VUgras{ٹُوٹی}
\textsuperscript{(70)} پھڑی \VUgras{دھار} \textsuperscript{(71)} نوں اگ دے وچکار
سادھدا اے۔

%%K t11-c2-04-1 p
4. --- رنگ منچ دے دوجے پاسے وڈی \VUgras{پوڑی} \textsuperscript{(72)} کھڑی کیتی گئی
اے۔ اوہ اگ بجھاؤ والیاں نوں اوہناں لاٹاں کول اپڑن دیندی اے جیہڑیاں کالے
\VUgras{دھوئیں} \textsuperscript{(75)} دے وَروَلیاں وچکار اُٹھدیاں نیں۔

%%K t11-tit-7 sub
\VUsaut{5.0mm}
\VUpk{10.2pt}{40}{سورمائی دے کم۔}
\VUsaut{3.4mm}

%%K t11-c2-05-1 p
5. --- \VUgras{اگ} \textsuperscript{(74)} \VUgras{رنگ منچ} \textsuperscript{(73)} دے
ڈیوڑھی کمرے وچ لگی، اوس \VUgras{اوپیرا} دی مشق دے ویلے جیدا پہلا \VUgras{کھیل} کل
ہوݨ والا سی۔ خوش قسمتی نال ہال وچ نِرے کجھ \VUgras{ویکھݨ والے} \textsuperscript{(76)}
سن۔ وچاریاں نے \VUgras{بالاخانے} \textsuperscript{(77)} وچ پناہ لئی، جِتھے سورمے
\VUgras{اگ بجھاؤ والے} \textsuperscript{(86)} پوڑی تے رسیاں نال اوہناں دی مدد نوں آ
رہے نیں۔

%%K t11-c2-06-1 p
6. --- \VUgras{ڈیوڑھی} \textsuperscript{(78)} دے تھلے، جِتھے
\VUgras{ناٹک دی سوچنا} \textsuperscript{(79)} کھیل دا اعلان کردی اے،
\VUgras{ساز والے} \textsuperscript{(81)} نسدے وکھائی دیندے نیں، جیہڑے
\VUgras{ساز منڈلی} \textsuperscript{(80)} دے نیں، اپݨے ساز چُکی:
\VUgras{وائلن} \textsuperscript{(81)}، \VUgras{بنسری} \textsuperscript{(82)}،
\VUgras{چیلو} \textsuperscript{(83)}، \VUgras{ٹرومبون} \textsuperscript{(84)}،
\VUgras{ڈھول} \textsuperscript{(85)}، وغیرہ۔ \VUgras{ساز سامان}
\textsuperscript{(87)} دا اک حصہ بچاؤݨ دا جتن کیتا جا رہیا اے۔

%%K t11-c2-07-1 p
7. --- \VUgras{اداکاراں} \textsuperscript{(89)} تے \VUgras{اداکارائیں}
\textsuperscript{(88)} نوں، گویاں تے گویاواں نوں، اپݨے منچ دیاں \VUgras{پوشاکاں}
\textsuperscript{(90)} وچ ای نسّݨا پیا اے۔ مکھ گویا اک \VUgras{صحافی}
\textsuperscript{(92)} نوں سمجھا رہیا اے پئی ایہہ کِویں ہویا۔ \VUgras{نامہ نگار}
پہلے ای پل توں افسراں نال دوڑیا آیا سی: \VUgras{منتظم} \textsuperscript{(91)}،
\VUgras{مئیر} \textsuperscript{(93)}، \VUgras{پولیس افسر} \textsuperscript{(94)} تے اک
\VUgras{جج} \textsuperscript{(95)}، جیہڑے پڑتال کرن گے پئی ذمے واری کیدی اے۔
\VUsaut{6.0mm}
\VUfilet{22mm}
